\section{The single-curve determinant and its central value}\label{sec:single}

Fix an odd prime \(p\).  In this section \(E/\Q\) is non-CM and has
analytic rank zero or one.  The purpose is to prove
\(X(E)\geq0\), where
\[
 X(E)=v_p(Q_E)-v_p\bigl(\#\Sha(E/\Q)\bigr).
\]
The rank and finiteness conclusions of
Theorem~\ref{classical:analytic-low-rank} are in force.
We first construct an integral element of a power-series
ring.  Its value at the trivial character will be computed separately:
by the dual exponential in rank zero, and by an exact logarithmic
comparison in rank one.  These are different assertions, and both are
needed for the inequality.

We use the inverse determinant and the local-condition conventions of
Section~\ref{sec:models}.  In particular, if a two-term acyclic complex
in degrees \(1,2\) has determinant of differential \(a\), then the
coordinate of its canonical rational trivialization has valuation
\(-v_p(a)\).  This sign will account for the local Euler factors below.

\subsection{Integral Siegel classes and their period coordinate}
\label{single:siegel}

Let \(N\) be the conductor of \(E\) and \(f\) its normalized
newform.  We construct the integral
classes needed for the single-curve comparison and determine their exact
dual-exponential coordinate.  The normalization follows the construction
in \cite[Section~4]{B}; the coefficient prime in that source is two, so
we verify the changes in the construction at an odd prime explicitly.

Let \(E_0\) be the optimal elliptic quotient of \(J_1(N)\), and let
\(C_{\rm cusp}\subset E_0(\overline{\Q})\) be generated by geometric
cusp differences.  Manin--Drinfeld makes this a finite Galois-stable
subgroup.  Replace the curve within its isogeny class by
\(E=E_0/C_{\rm cusp}\); Proposition~\ref{classical:isogeny} preserves
\(X(E)\).  The Abel map, followed by this quotient, gives
\(\pi:(X_1(N),D_N)\longrightarrow(E,O)\), where \(D_N\) is the cusp set;
the map is independent of the base cusp used for the Abel map.
Put \(T=H^1(\overline E,\Z_p)(1)=T_pE\), using the principal pairing,
and \(V=T\otimes\Q_p\).
For \(A=p^k\), \(k\geq0\), put \(L_0=NA\) and let
\(\Pi:(X_1(L_0),D_{L_0})\to(E,O)\) be the \(A\)-degeneracy followed
by \(\pi\).  At the mark \(1/L_0\) the degeneracy is
\(\tau\mapsto A\tau\).  Thus, for a minimal differential \(\omega\),
\[
 \Pi^*\omega=c_FF\,\frac{dq}{q},\qquad
 F(\tau)=Af(A\tau),\qquad q=e^{2\pi i\tau},\quad c_F\in\Q^\times.
\]
Let \(a\) be the primitive generator of
\(H_1(E(\C),\Z)^{G_{\R}}\) oriented by
\(\Omega_0=\int_a\omega>0\).  For \(\xi\in\mathrm{SL}_2(\Z)\),
define \(\delta_\xi\in H^1(E(\C),\Z)\) as the Poincare dual of
\(\Pi_*\xi\{0,\infty\}\), with convention
\(\int_E\omega\wedge\delta_\xi=\int_{\Pi_*\xi\{0,\infty\}}\omega\).

\begin{lemma}[The relative lattice]\label{single:relative-lattice}
Relative pullback and Poincare duality give an integral Galois map
\[
 H^1(\overline{Y_1(L_0)},\Z_p)(1)\longrightarrow T.
\]
For \(A=1\), an integral combination \(\delta=\sum_\xi n_\xi\delta_\xi\)
of finitely many symbols satisfies \(\delta(a)=1\).
\end{lemma}
\begin{proof}
Write \(E_0(\C)=\C/\Lambda_0\).  The connected optimal kernel makes
the map on integral homology onto \(\Lambda_0\).  Quotienting by
\(C_{\rm cusp}\) replaces \(\Lambda_0\) by the lattice generated by it
and lifts of the cusp images.  Relative paths between cusps supply
exactly these additional generators.  Consequently
\(H_1(X_1(N)(\C),D_N;\Z)\to H_1(E(\C),\Z)\) is onto.
The Farey triangulation generates the relative group by the Manin edges
\(\xi\{0,\infty\}\).  The invariant sublattice is saturated, so \(a\)
is primitive; unimodularity of the intersection form now gives the
asserted integral combination.  Pullback by the map of pairs is integral,
and relative Poincare duality gives the displayed etale map.  This is
the Betti and relative-cohomology argument of
\cite[Section~4, ``The integral relative lattice'']{B}, before any choice
of coefficient prime; no Hecke projector is involved.
\end{proof}

We use two level patterns, always with \(v_p(M)\geq6\).
In square order, \(n_1=n_2=M\), \(L_0\mid M\), and the row
\(l=(l_1,l_2)\) is \((0,1)\xi\).
In rectangular order put \(b=N/p^{v_p(N)}\), require
\((b,M)=1\) and \(p^{v_p(L_0)}\mid M\), and take
\[
 n_1=b_1M,\qquad n_2=bM,\qquad l=(Ns,1)=(0,1)\xi,
\]
where \(b_1\mid b\) is a product of complete prime powers dividing \(b\).
The unfilled pattern means \(b_1=1\).
The curve \(Y(n_1,n_2)\) parametrizes elliptic curves with independent
points \(P,Q\) of orders \(n_1,n_2\); independence means that they
embed \(\Z/n_1\oplus\Z/n_2\).  The maps to the marked curve and the
cyclotomic field are specified by
\begin{equation}\label{single:mark-root}
 C=\frac{l_1n_1}{L_0}P+\frac{l_2n_2}{L_0}Q,\qquad
 \eta_M=e_M((n_1/M)P,(n_2/M)Q).
\end{equation}
The coefficients of \(C\) are integers in both patterns; \(C\) has
order \(L_0\), and \(\eta_M\) is primitive of order \(M\).
Fix compatible roots in the algebraic, complex, and \(p\)-adic
realizations, with complex root \(e^{2\pi i/m}\) at order \(m\).

Choose \(c,d>1\) prime to \(6pn_2\), with
\(c,d\equiv1\pmod{L_0}\).  On the universal elliptic curve \(\mathcal E\)
let \({}_c\theta_{\mathcal E}\) be Kato's normalized function of divisor
\(c^2(0)-\mathcal E[c]\), and put
\({}_cg_P=P^*{}_c\theta_{\mathcal E}\), \({}_dg_Q=Q^*{}_d\theta_{\mathcal E}\).
Its norm under \([s]\) is itself if \((s,c)=1\)
\cite[Sections~1--2]{Kato}.
At precision \(p^j\), replace the orders by \(p^jn_1,p^jn_2\), and
let \(P',Q'\) range over their independent division lifts of \(P,Q\).
Take the unnormalized trace of the Kummer cup
\begin{equation}\label{single:K1}
 ({}_cg_{P'})\cup({}_dg_{Q'})\otimes
 e_{p^j}(n_1P',n_2Q')^\vee.
\end{equation}
Here the dual basis sends the displayed primitive root to \(1\), so the
coefficient is \(\Z/p^j(1)\).  Apply the arithmetic edge map, trace
along \eqref{single:mark-root}, and push to \(T/p^j\) by
Lemma~\ref{single:relative-lattice}.  Passage to the inverse limit
defines a class in \(H^1(G_{\Q,S},
T\otimes\Z_p[\Gal(\Q(\mu_M)/\Q)])\), for \(S\) containing the
primes of \(pn_2\) and infinity.
For a finite character \(\nu\) of this Galois group, send the basis
vector at \(\sigma\) to \(\nu(\sigma)\); write \(z_\nu\) for the result
with coefficients in \(T\otimes\Z_p[\nu]\).
All traces, including this character push, are unnormalized.

\begin{proposition}[Integrality and traces]\label{single:siegel-traces}
The classes just constructed are integral away from \(pn_2\infty\)
and compatible with division traces and with corestriction as the
\(p\)-part of \(M\) grows with \(A,c,d,\xi\) fixed.
If \(\ell\nmid cdn_2L_0\) is a fresh prime and is adjoined to both
orders and the root index, the rational norm multiplier is
\begin{equation}\label{single:good-norm}
 1-\frac{a_\ell(f)[\ell]}{\ell}+\frac{[\ell]^2}{\ell},
\end{equation}
where \([\ell]\) sends the determinant root to its \(\ell\)-th power.
For \(q_0^a\parallel b\), \(q_0\nmid b_1\), filling \(q_0^a\)
into the first row gives, at a finite even character \(\nu\),
\begin{equation}\label{single:bad-row}
 z_{\rm filled}=(1-a_{q_0}(f)\nu(q_0)/q_0)z_{\rm before}.
\end{equation}
The trace maps themselves are integral in both assertions.
\end{proposition}
\begin{proof}
Further \(p\)-division gives a Cartesian product of the two lift sets:
independence is automatic since the old tuple is already a basis at
\(p\).  At the old precision the inverse-root twists agree.  Base
change and the projection formula turn the trace of the cup into
the cup of the traces, and distribution recovers the old class.
The mark and lower root also agree when the base \(p\)-level grows.
The functions and their inverses are integral away from the level.
Indeed horizontal valuations follow from their divisors; on a smooth
full-level model every geometric fiber component meets a cusp, where
the Tate leading terms are signs, roots of unity, and powers of
\(1-\zeta\) for nontrivial level roots.  These are units away from
the level.  This also checks quotients of level after pullback.
The relative open curve has lisse geometric cohomology away from this
support, vanishing above degree one.  Its arithmetic edge map lands in
\(H^1(G_{\Q,S},H^1(\overline{Y(n_1,n_2)},\Z/p^j)(1))\).
Edge maps commute with traces and maps of pairs.  Finiteness in degree
zero gives the Mittag--Leffler property needed for compact passage;
the subsequent lattice and coefficient maps are integral.

At \(\ell\), start with the Cartesian division sets, subtract the
tuples whose \(\ell\)-parts lie in a common line \(I\subset\mathcal E[\ell]\),
and add back \(\ell\) times the double-prime-to-\(\ell\) lift.  The zero
pair was subtracted \(\ell+1\) times, which explains this multiplicity.
Distribution on \(\mathcal E/I\) identifies the line sums with the
good Hecke correspondence.  Quotient isogenies raise the Weil roots
to \(\ell\), giving \([\ell]\) and twist scalar \(\ell^{-1}\);
the relative elliptic push has Hecke eigenvalue \(a_\ell(f)\).
This last assertion holds for relative maps because pullback has zero
endpoint-function coordinates and the equality of eigenforms also
holds on path integrals.  The added-back multiplication isogeny raises
the roots to \(\ell^2\); its scalar \(\ell^{-2}\), multiplied by
\(\ell\), gives the last term of \eqref{single:good-norm}.

For row filling, fix \(Q\).  A discarded \(q_0^a\)-division lift of
\(P\) has its \([q_0^{a-1}]\)-image in the order-\(q_0\) line \(I_Q\).
For \(\rho:\mathcal E\to\mathcal E/I_Q\), these lifts are precisely
the preimages under \(\rho[q_0^{a-1}]\) of the prime-to-\(q_0\) point
\(\widetilde P\) satisfying \(\rho^\vee\widetilde P=P\).
Express the other unit by distribution over
\(\rho^\vee\widetilde Q=Q\).  The points \(\widetilde Q\) have order
\(n_2\), and their order-\(q_0\) lines differ from \(\ker\rho^\vee\).
In reverse, the discarded sum is indexed by all lines other than the
line of the source mark, since \(l=(Ns,1)\).  Its branches at mark
\(1/L_0\) are \(\tau\mapsto(\tau+i)/q_0\), \(i\bmod q_0\).
Thus it is \(U_{q_0}\), acting on \(Af(A\tau)\) by \(a_{q_0}(f)\)
because \((A,q_0)=1\).  Root translation and inverse twist contribute
\(\nu(q_0)/q_0\).  This proves \eqref{single:bad-row}, including
\(q_0=2\).  The argument reproduces the geometric correspondences of
\cite[Section~4, ``Good-prime norms and rectangular filling'']{B};
every denominator displayed here is a unit at \(p\).
\end{proof}

We next fix the rational differential coordinate, retaining the absolute
normalization needed for divisibility.  For a nonzero torsion label
\(P=[x\tau+y]\), let \(E_P\) be the algebraic weight-one form
\[
 E_P=-\frac1{2\pi i}\left.
 \sum_{(m,n)\in(x,y)+\Z^2}
 \frac{(\operatorname{Im}\tau)^w}
 {(m\tau+n)|m\tau+n|^{2w}}\right|_{w=0},
\]
where the value is obtained by regular continuation.

\begin{lemma}[Full-level calibration]\label{single:kato-calibration}
In square order, before projection to the elliptic curve, the dual
exponential of \eqref{single:K1} is
\begin{equation}\label{single:full-differential}
 \pm M^{-2}(c^2E_P-cE_{cP})(d^2E_Q-dE_{dQ})\,\frac{dq}{q}.
\end{equation}
The sign is determined by cup, pairing, and orientation conventions
and is common to the determinant components.
\end{lemma}
\begin{proof}
Initially impose \(c,d\equiv1\pmod M\).  Apply
\cite[Theorem~9.5]{Kato} with \(k=2\), \(r=r'=1\), and equal level
orders \(M,M\).  The prime-support containment is automatic, the
range conditions are satisfied, and \(p\mid M\) selects the branch
without an additional Euler operator.  The symmetric power has degree
zero.  Its negative Tate moment is exactly the inverse determinant
\(e_{p^j}(MP',MQ')^\vee\) on the basis coset of division tuples.
The weight-one shifted and Fourier-dual forms agree under the fixed
conventions \cite[Proposition~3.11(1) and Section~4.2]{Kato}.
The theorem therefore gives proportionality to the two brackets in
\eqref{single:full-differential}.  We determine the scalar by residues.

Put \(\Theta(t,q)=(1-t)\prod_{h\geq1}(1-q^ht)(1-q^h/t)\).
The normalized Tate function is
\[
 {}_c\theta(t)=(-1)^{(c-1)/2}q^{(c^2-1)/12}t^{-c(c-1)/2}
                 \frac{\Theta(t,q)^{c^2}}{\Theta(t^c,q)}.
\]
The relation \(\Theta(qt,q)=-t^{-1}\Theta(t,q)\) verifies periodicity
and the divisor; multiplication over \([s]\)-preimages verifies norm
invariance for \((s,c)=1\).  These norms fix the constant uniquely:
the norms under \([2]\) and \([3]\) force its third and eighth powers
to be one.  For \(0<x<1\), \(cx\notin\Z\), put
\[
 J_c(x)=c\lfloor cx\rfloor-\frac{c(c-1)}2
       =c^2B_1(x)-cB_1(\{cx\}),\qquad B_1(x)=x-\tfrac12.
\]
At \(t=q^x\zeta\) the leading coefficient is
\(\pm\zeta^{J_c(x)}\), and the order is
\(xJ_c(x)+(c^2-1)/12-\lfloor cx\rfloor(\lfloor cx\rfloor+1)/2\).
The last two terms are integers, since \((c,6)=1\).

On a determinant component choose the cusp
\(P=q^{1/M}\), \(Q=q^{1/M}\zeta_M^u\), with \((u,M)=1\),
over \(K_0=\Q_p(\mu_M)\).  Set \(M_j=p^jM\) and
\(K_j=K_0(\mu_{M_j})\).  Each pair
\(a_1,a_2\in[0,M_j)\cap(1+M\Z)\) gives a cusp orbit represented by
\(P'=q^{a_1/M_j}\), \(Q'=q^{a_2/M_j}\zeta_{M_j}^u\).
Translation by the old width eliminates the first root index; determinant
primitivity makes cyclotomic Galois act freely on the remaining indices.
There are \(p^{2j}\) such orbits, residue degree \(p^j\), and width
ratio \(M_j/M=p^j\), accounting for \(p^{4j}\) division tuples.
Prime-to-\(p\) roots contribute only fixed data.
With \(x=a_1/M_j\), \(y=a_2/M_j\), the tame symbol has root exponent
\(h\equiv\pm uJ_d(y)a_1J_c(x)\pmod{p^j}\).
All signs have zero Kummer class because \(p\) is odd.  The inverse
twist uses \(\zeta_{p^j}^{\pm a_1u}\), so the twisted residue on
\(G_{K_j}\) is
\(\pm J_c(x)J_d(y)(\chi_{\rm cy}-1)/M_j\pmod{p^j}\).
Abelian transfer and the exponential series give
\[
 \frac{\chi_{\rm cy}^{p^j}-1}{p^jM}
       \equiv\frac{\log\chi_{\rm cy}}M\pmod{p^j}.
\]
Indeed \(\log\chi_{\rm cy}\in p^{v_p(M)}\Z_p\), and every term
of degree at least two has the required valuation.  Residue trace
sums field transfers without a ramification multiplier.  Bernoulli
distribution gives \(\sum_{a_i}J_c(a_i/M_j)=J_c(1/M)\), and likewise
for \(d\); hence the limiting residue is
\(\pm J_c(1/M)J_d(1/M)\log\chi_{\rm cy}/M\).
The dual exponential sends \(\log\chi_{\rm cy}\) to \(\pm1\): pair
with Kummer classes of exponential units and use local reciprocity,
with field trace in the de Rham pairing.  The constant term of \(E_P\)
is \(-B_1(x)\), so the two brackets have nonzero constant terms
\(-J_c(1/M)\), \(-J_d(1/M)\) for the initial choices.  Since
\(dq/q\) has residue \(M\), the scalar is \(\pm M^{-2}\).

For arbitrary permitted smoothings use
\((c_0^2-[c_0]^*){}_c\theta=(c^2-[c]^*){}_{c_0}\theta\).
Changing a row by \(c\) changes its inverse-root basis by \(c^{-1}\),
so the induced differential operator is \(c^2-c\,\mathrm{mult}_c^*\).
Choose \(c_0\equiv1\pmod M\), cancel \(c_0^2-c_0\) rationally,
and repeat on the other row.  This cancellation asserts a rational
identity only; it introduces no claim of integral divisibility.
\end{proof}

\begin{proposition}[The period coordinate]\label{single:period-law}
In square or unfilled rectangular order, for every finite even
character \(\nu\) of conductor dividing \(M\), the coordinate of
\(\exp^*(z_\nu)\) in \(\omega\), in compatible field realizations, is
\begin{equation}\tag{2}\label{eq:source-2}
 \pm\Delta_\nu\frac{L^{(M)}(f,\nu,1)}{\Omega_0}\delta_\xi(a),\qquad
 \Delta_\nu=(c^2-c\nu(c)^{-1})(d^2-d\nu(d)^{-1}).
\end{equation}
The superscript omits all primes dividing \(M\), and the sign is
common across symbols.  In square order with \(A=1\), an integral
combination has \(\delta(a)=1\).  In rectangular order, some choice
of \(A=p^k\) and \(\xi\) with lower row \((Ns,1)\) has
\(\delta_\xi(a)\ne0\).
\end{proposition}
\begin{proof}
Apply the coefficient-prime-independent analytic trace identity of
\cite[Section~4, equation~(K5) and proof of ``The period formula'']{B}
to \eqref{single:full-differential}.  Its inputs are square order,
\(L_0\mid M\), mark \(1/L_0\), even \(\nu\), and
\(c,d\equiv1\pmod{L_0}\), all imposed here.  The two rescaled
weight-one lattice sums cancel \(M^{-2}\).  Row reindexing preserves
the mark and gives \(\Delta_\nu\); Hecke unfolding has indices prime
to \(M\) and gives \(L^{(M)}(f,\nu,1)\).
The relative de Rham push has zero cusp-function coordinates, and its
pairing with \(\overline\omega\) is
\(\pm2i\operatorname{Im}(\int_{\Pi_*\xi\{0,\infty\}}\omega)
\Delta_\nu L^{(M)}(f,\nu,1)\).  Writing
\(\delta_\xi=s_0\omega+\overline{s_0}\overline\omega\) gives
\(s_0+\overline{s_0}=\delta_\xi(a)/\Omega_0\), which proves (2).
These are algebraic differentials and algebraic trace and push maps;
the same identity therefore gives the specified \(p\)-adic coordinate.
Local duality uses field trace, with no character-orbit average.
For rectangular order fill every missing row prime.  Each factor
\(1-a_{q_0}(f)\nu(q_0)/q_0\) is nonzero since \(|a_{q_0}(f)|\leq1\),
and is precisely the additional omitted Euler factor at \(q_0\).
Canceling these factors rationally proves (2) before filling.

The square assertion follows from Lemma~\ref{single:relative-lattice}.
For the rectangular assertion take \(\xi\) with upper row \((1,0)\).
The Fricke relation sends the degenerated path from \(0\) to
\(p^k/(Ns)\) to the path from \(\infty\) to \(-s/p^k\), up to sign.
Rohrlich's theorem in \cite[Theorem~13.5(2)]{Kato}, with prime support
\(\{p\}\), gives a primitive odd character of sufficiently high
\(p\)-power conductor with nonzero central twist.  Odd characters
exist even at \(p=3\), by adjoining a fixed odd tame character to
deep wild characters.  Primitive Gauss--Mellin then makes the weighted
sum of these path integrals nonzero.  If they were all real, complex
conjugation would identify the integrals at \(s\) and \(-s\), and
their sum against an odd character would vanish.  Some imaginary part,
and therefore the corresponding \(\delta_\xi(a)\), is nonzero.
\end{proof}

Finally, rational norm identities can be made integral after one common
nonzero multiplier, uniformly over the abelian fields used here.
By Serre's open-image theorem \cite{Serre}, the commutator image on
\(T\) contains a fixed deep special-linear congruence subgroup.
Its invariants in \(T\otimes\Q_p/\Z_p\) have bounded exponent.
The exact sequence for \(T\subset T\otimes\Q_p\) consequently bounds
the torsion exponent of \(H^1(F,T)\) uniformly for these abelian fields
\(F\).  Multiplying a rationally zero norm error by this bound kills it.

\subsection{The determinant to be made integral}

Let \(S\) contain \(p,\infty\) and every bad prime of \(E\), and fix
the smoothing integers before choosing moving primes.  Their prime
divisors are outside the finite support used for the square classes.
Take the rational auxiliary tuple of Section~\ref{sec:vertical}, with
one tame variable for each moving prime and a genuine cyclotomic
variable \(t\).  Put
\[
 R=\Z_p[[t,u_1,\ldots,u_k]],\qquad M_R=T\otimes_{\Z_p}R(\chi).
\]
Here \(\chi\) takes the chosen generators to \(1+t\) and
\(1+u_j\).  The complex \(G\) is the limiting unrestricted global
complex, with the moving primes allowed in its support.  We retain
the finite-stage complexes with their actual cyclotomic tower, as
well as their reductions modulo every fixed Artin precision.  Let
\(z\in H^1(G)\) be the smoothed square class with symbol evaluation
one and degeneracy \(A=1\).  At each stage choose the square order
\(M\) with prime support exactly \(S_f\) together with the moving
primes, and divisible by the required fixed orders.  Thus its
allowed ramification support and its omitted Euler factors agree.

For a fixed \(q\in S_f\setminus\{p\}\), define
\[
 D_q=\det\left(1-\frac{\chi(\Frob_q)\rho_T(\Frob_q)}q
                         \,\middle|\,V_{I_q}\right).
\]
The subscript denotes inertia coinvariants.  Let \(P_q\) denote the
inverse local \(L\)-polynomial at \(1\), with its unramified scalar
argument in the orientation of the period formula.  The two scalar
arguments may be inverse to one another; at the trivial character
they agree.  The polynomials are respectively quadratic at good
reduction, signed linear at multiplicative reduction, and one at
additive reduction.  Thus they are integral, including when \(p\)
divides a Tamagawa number.  Moreover
\[
 D_q(0)=P_q(0)=L_q(E,1)^{-1}\ne0.
\]
We also use the smoothing series \(\Delta\) of
\eqref{eq:source-2}.  Each of \(D_q,P_q,\Delta\) remains nonzero
modulo \((p,u_1,\ldots,u_k)\): the cyclotomic exponents of the
positive integers \(q,c,d\) are nonzero.  Consequently these factors
have no vertical divisor and are regular in \(t\).

\begin{proposition}[Integral normalized coordinate]\label{single:integral}
The generic cohomology of \(G\) is a line in degree one, and \(z\)
is nonzero on that line.  If \(U_0\in\Frac R\) is its coordinate in
the inverse determinant relative to an integral generator, then
\begin{equation}
 U=\frac{U_0}{\Delta}
       \prod_{q\in S_f\setminus\{p\}}\frac{D_q}{P_q}
       \ \in R\setminus\{0\}.
 \tag{3}\label{eq:source-3}
\end{equation}
\end{proposition}

The argument adapts the horizontal comparisons and division method
of \cite[Section~5, ``Horizontal line switches'' and
``Removal of the fixed divisors'']{B}.  The proof has three parts.
The concentration assertion
\eqref{eq:source-1}(a) treats the divisor \((p)\).  At a
characteristic-zero divisor, auxiliary Euler-system derivatives
compare the coordinate with a coordinate on a free line.  Finally,
rectangular classes remove the divisors of the smoothing and fixed
Euler factors.  The last step is necessary: integrality of the
smoothed class alone does not permit division by \(\Delta\).

\subsection{Horizontal line comparisons}

We use the following three settings.  In the first, the coefficient
ring is the localization of \(R\) at a height-one prime avoiding
\(p\Delta\prod_{q\ne p}D_qP_q\).  In the second, all tame
parameters have been removed and we use the fraction field of the
fixed-support cyclotomic ring.  In the third, at an actual stage,
every tame parameter is specialized to a nontrivial finite
\(p\)-power character \(\eta_j\), while \(t\) is retained.  Write
\(\eta=\prod_j\eta_j\) and
\(\Lambda_\eta=\Z_p[\boldsymbol\eta][[t]]\).  The third setting
uses its generic field or a height-one localization not above \(p\),
and imposes unramified conditions at the old non-\(p\) places.
We require that a fixed Frobenius lift at \(p\), trivial in the
cyclotomic quotient, already kill the dual invariants after twisting
by \(\eta\).  The nonzero exponent vector at \(p\) supplies such
specializations.

\begin{lemma}[Line comparison for a smoothed system]\label{single:line}
In any of these three settings, suppose that the chosen square or
rectangular system class \(z_{\rm sys}\) is generically nonzero.
Then the global complex has a single generic cohomology line in
degree one.  At each of the indicated DVRs,
\begin{equation}
                  d(G_{\rm test},z_{\rm sys})\geq0.
                  \tag{4}\label{eq:source-4}
\end{equation}
In the third setting, \(G_{\rm test}\) denotes the complex with the
unramified old local conditions.
\end{lemma}

\begin{proof}
We verify the local hypotheses before constructing derivatives.
In the first setting, dual local invariants vanish at a fixed
non-\(p\) place by the exclusion of \(D_q\).  At a moving place,
invariance under inertia first forces \(u_j=0\); the remaining
Frobenius equation is nonzero on that divisor because the limiting
Frobenius has determinant one and trace different from two.
At \(p\), choose an inertia element generating the cyclotomic
quotient and a Frobenius lift trivial in that quotient.  The first
invariance equation is regular in \(t\); the second is independent
of \(t\) and nonzero, by the fixed-place exponent condition.  They
have no common height-one component.  In the second setting,
nonconstant cyclotomic action gives the same vanishing.  In the
third, the imposed condition at \(p\) gives it directly.

At an old non-\(p\) place in the third setting, inertia acts through
a constant representation.  The unramified condition is residue
cochains on inertia invariants.  Its quotient is residue cochains
on inertia coinvariants with twist \((-1)\), in degrees \(1,2\).
Its differential has determinant
\[
 D_{q,\eta}=\det\left(1-\frac{\chi_t(\Frob_q)
                     \rho_{V\otimes\eta}(\Frob_q)}q
                         \,\middle|\,(V\otimes\eta)_{I_q}\right).
\]
This is nonzero and equals one at a moving conductor where \(\eta\)
is ramified.  The quotient therefore has zero \(H^1\) over the
DVR itself.  This assertion is made before passing to its residue
field, where the determinant may vanish.  The conditions include
all \(H^0\) and omit \(H^2\).  Unramified cup products factor through
residue cohomology, so are zero; local duality and the Euler
characteristic give exact complementarity, including at the
residue field.  These are precisely the complex conditions needed
by Lemma~\ref{models:line-switch}.

Global invariants vanish by the non-CM open-image argument
\cite{Serre}.  We spell out its use on the limiting diagrams.  In
the product of the stage absolute groups, take the joint kernel of
the Tate action and the abelian and radical data used in the
characters.  The quotient contains arbitrary sequences from a
fixed deep principal subgroup of \(\operatorname{SL}_2(\Z_p)\),
acting trivially on the latter data.  Indeed those data have
bounded derived length, whereas repeated Lie brackets retain
\(\mathfrak{sl}_2\).  The evaluation and diagonal-conjugation
argument of Lemma~\ref{models:evaluations} shows that restriction to
this kernel detects a nonzero degree-one class.  Its evaluations
span the Tate plane by irreducibility.  Both the primal and the
opposite-condition spaces inject into unrestricted global
cohomology, by the local invariant checks above.  Poitou--Tate and
the global Euler characteristic now give amplitude \(1,2\) and
\[
 h=1+h^*,
\]
where \(h,h^*\) are their degree-one dimensions.  This remains
true after any of the line changes below.

Choose a fresh sequence of good primes \(\ell\to1\) in
\(\Z_p^\times\), split to increasing precision on all old scalar
data, whose Tate Frobenius tends to
\[
 \tau=\begin{pmatrix}1&e\\0&1\end{pmatrix},\qquad e\ne0.
\]
The local degree-one cohomology over the testing DVR splits into
the finite line \(T/(\tau-1)T\), measured by Frobenius, and the
singular line \(\ker(\tau-1)\), measured by tame inertia.
The notation here means scalar extension to the testing ring;
\(e\) is a unit there because \(p\) has been inverted.  The
conjugating power in the inertia relation tends to one.  The
finite line is isotropic by inflation, and a pure singular cocycle
has values in the Weil-isotropic invariant line.  The mixed pairing
is perfect.  Taking all \(H^0\), the respective line in \(H^1\),
and no \(H^2\) therefore gives the exact complementary conditions.

A nonzero primal class and a nonzero dual class can both be made
nonzero on the finite line.  On the joint kernel their evaluations
are additive and span the plane.  Each prohibited zero is a proper
affine condition; finitely many such conditions can be avoided in
characteristic zero.  Chebotarev realizes the resulting finite
tables at increasing precision.

The order of the further choices is relevant.  At each stage fix
the cyclotomic base \(K_0=\Q(\mu_M)\) first.  Perturb the desired
unipotent element within the available deep special linear
subgroup so that
\(v_p\det(1-\rho_T(\Frob_\ell))\) is large but finite.  Next
choose \(A_i\) with
\[
 A_i-v_p\det(1-\rho_T(\Frob_\ell))\longrightarrow\infty,
 \qquad \ell\equiv1\pmod{p^{A_i}}.
\]
Require exact splitting in \(K_0\), all earlier derivative fields,
and the \(p^{A_i}\)-th roots of all earlier derivative primes.
These conditions make the new degree-\(p^{A_i}\) conductor-\(\ell\)
field and the earlier derivative fields mutually split at their
conductors.  They are compatible with the finite Tate and
evaluation conditions by the same open-image argument.  A
Frobenius lift trivial in the new derivative field is obtained by
an inertia adjustment.

For a set \(I\) of derivative primes, norm the system class to the
compositum \(K_I\) of their degree-\(p^{A_i}\) fields with \(K_0\),
choose inertia representatives \(\sigma_\ell\) generating their
own cyclic factors and trivial on the other factors, and apply
\[
 D_I^{\rm der}=\prod_{\ell\in I}
                 \sum_{j=1}^{m_\ell-1}j\sigma_\ell^j,
 \qquad m_\ell=p^{A_i}.
\]
The good Euler norm at \(\ell\) is
\[
 p_\ell=\ell^{-1}\det(1-\rho_T(\Frob_\ell)).
\]
There is one fixed nonzero \(p\)-power \(L_*\) which clears all
torsion ambiguities and descent obstructions used in a finite
switching argument.  To see uniformity, the commutator image for
a non-CM curve bounds the exponent of \((V/T)^{G_F}\) over all
abelian fields \(F\) in question.  This bounds \(H^1(T)\)-torsion,
and also the invariants of \(T/p^nT\) that occur in
inflation--restriction.  The obstruction and ambiguity groups
are killed by the same bounded exponent.  The identity
\[
 (\sigma_\ell-1)D_\ell^{\rm der}=m_\ell-N_\ell,
 \qquad N_\ell=\sum_{j=0}^{m_\ell-1}\sigma_\ell^j,
\]
then makes the derivatives invariant modulo \(p^n\), for
\(n\to\infty\) sufficiently slowly.  Descend to \(K_0\), use
unnormalized character corestriction, and pass to the marked
models.  We obtain integral classes \(Z_I\), with
\(Z_\varnothing=L_*z_{\rm sys}\).

Here is the exact local derivative comparison.  Put \(m=m_\ell\).
Let \(b,b_\ell\)
be the lower and upper integral cocycles after all other
derivatives, and put \(F=\rho_T(\Frob_\ell)\).  They vanish on
local inertia: good reduction and the Weil bound exclude the
eigenvalue \(\ell\) in the relevant inertia cohomology.  Put
\(x=b(\Frob_\ell)\), \(x_\ell=b_\ell(\Frob_\ell)\).  The norm
identity gives a coboundary \(dy\) with
\[
 \sum_{j=0}^{m-1}\sigma_\ell^jb_\ell=p_\ell b+dy,
 \qquad (F-1)y=mx_\ell-p_\ell x.
\]
After adjusting the descended cocycle, its restriction is exactly
the derivative of \(b_\ell\).  Its finite value tends to zero
because the derivative coefficient sum is \(m(m-1)/2\).  Applying
\(\sigma_\ell-1\) shows that its singular value is \(-y\), up
to the uniformly bounded loss already cleared.  Finally
\[
 y=m(F-1)^{-1}x_\ell-\ell^{-1}\operatorname{adj}(F-1)x.
\]
The first term tends to zero by the chosen margin on \(A_i\).
A lower coboundary \((F-1)w\) disappears after applying the
adjugate in the limit.  The calculation holds on every translate
before character corestriction.  Therefore
\begin{equation}\label{single:finite-singular}
 f(Z_I)=0,\qquad
 s(Z_I)=\operatorname{adj}(\tau-1)f(Z_{I\setminus\{\ell\}}).
\end{equation}
The map between the two lines is multiplication by \(\pm e\),
and hence is a unit comparison at the testing DVR.  The classes
lift to the complexes with the new conditions; the inclusion of
\(H^0\) and the vanishing of the old singular quotients' \(H^1\)
give the claimed lifts over the DVR itself.

If \(h>1\) generically, choose the two nonzero evaluations above.
Equation~\eqref{single:finite-singular} makes the derived primal
class nonzero on the singular line.  Its pairing with the old
dual class is then nonzero at this prime and zero at all other
places, contradicting reciprocity.  Thus \(h=1\) generically.
At a DVR where the fiber has \(h>1\), use the primitive reduction
of the generic line and a nonzero fiber dual class.  The exact
line-switch Lemma~\ref{models:line-switch} reduces \(h\) by one
and preserves the determinant valuation, because the comparison in
\eqref{single:finite-singular} is a unit.  After finitely many
switches the complex is a free line in degree one and its class
is integral.  The factor \(L_*\) is a unit at this DVR.  Reading
the equalities backwards proves \eqref{eq:source-4}.
\end{proof}

\subsection{Removing the smoothing and Euler divisors}

\begin{proof}[Proof of Proposition~\ref{single:integral}]
At \((p)\), \eqref{eq:source-1}(a) identifies \(G\) with a free
line in degree one.  The class is integral and the normalizing
factors are units there.  This proves nonnegative order at the
vertical divisor once generic nonvanishing is known.

For that nonvanishing, use the second setting in
Lemma~\ref{single:line}.  Rohrlich's finite-support twist
nonvanishing, in the form \cite[Theorem 13.5(2)]{Kato}, applies to
the fixed weight-two newform and characters with conductor
supported at \(p\).  Infinitely many sufficiently ramified even
characters have nonzero central value.  The period formula
\eqref{eq:source-2} therefore makes the fixed-support square class
generically nonzero: a torsion class over the one-variable
power-series ring would vanish at all but finitely many such
characters.  Lemma~\ref{single:line} gives one generic line.
Setting all \(u_j=0\) in the full diagram adds acyclic singular
quotients, since the moving Frobenius limits have
\(\det(1-g_j)\ne0\), where \(g_j\) is the limiting Frobenius
at the \(j\)-th moving prime.  The good norm relations multiply the class
by the nonzero constants \(2-\Tr(g_j)\).  Thus the full diagram
also has a nonzero class on a single generic line.  The first
setting of the lemma now treats every horizontal divisor outside
the fixed factors.

It remains to remove the fixed factors.  Express \(U=g_0/f_0\)
using integral determinant pivot formulas, with \(f_0\ne0\), and
retain integral stage expressions \(g_{0,i},f_{0,i}\).  Specialize
the tame variables to finite nontrivial values
\(u_j=\eta_j-1\), avoiding identically zero denominators and
the prohibited Frobenius values at \(p\).  These specializations
still detect bounded series: repeated one-variable Weierstrass
preparation permits avoidance of any finite collection of
nonzero conditions.  The conditions also hold at sufficiently
precise original stages.

Fix such a stage and a horizontal prime of \(\Lambda_\eta\).
Use a rectangular system with unfilled first row, with
\(b=N_{p'}\), and \(M\) supported only at \(p\) and the active
tame conductors.  Choose a degeneracy \(A=p^k\) and a symbol
with nonzero evaluation.  Smoothing integers \(c',d'\) can be
chosen so that \(\Delta'\) is a unit at this prime.  Indeed
preparation realizes its \(t\)-value as an algebraic point of
the open unit disk.  Use CRT to choose integers tending to one
\(p\)-adically, congruent to one on the required fixed symbol
orders, and with a prescribed nontrivial \(\eta\)-value at one
active conductor and trivial values at the others.  All
coprimalities can be maintained.  At the selected \(t\)-value,
the cyclotomic factor tends to one whereas the chosen constant
weight is nontrivial, so the smoothing factor is eventually
nonzero.  Fix one such choice while applying the line comparison.

Call the resulting class \(z'\).  It has primitive Euler factors
away from \(p\): the only other omitted primes are ramified
conductors of \(\eta\), where the twisted local factor is one.
The same nonvanishing theorem makes \(z'\) generically nonzero.
Lemma~\ref{single:line} gives its integral determinant order with
unramified old conditions.  On changing those conditions to full,
the order changes by
\(-\sum_{q\ne p}v(D_{q,\eta})\).

In the common generic full cohomology line the ratio of the
square and rectangular classes, apart from their nonzero
constant symbol ratio, is
\begin{equation}\label{single:class-ratio}
          \frac{\Delta\prod_{q\in S_f\setminus\{p\}}P_q}
               {\Delta'}.
\end{equation}
To verify the identity rather than merely its character values,
evaluate the two classes at the infinitely many sufficiently
ramified cyclotomic characters with nonzero central value.
Formula~\eqref{eq:source-2} gives exactly the displayed ratio;
omissions at active conductors contribute one.  Clearing the
generic denominators and applying one-variable preparation proves
the identity.  Squaring first gives the same valuation conclusion
if the common orientation sign has not been fixed.  Consequently
the tested quotient \(g_{0,i}/f_{0,i}\) has no horizontal pole,
that is,
\[
 g_{0,i}(t,\boldsymbol\eta-1)/f_{0,i}(t,\boldsymbol\eta-1)
                       \in\Lambda_\eta[1/p].
\]
There is no assertion of a uniform denominator for these stage
quotients.

Every remaining possible pole of \(U\) divides a factor regular
in \(t\).  For such a factor \(H\), take a distinguished
representative and divide the integral numerator by \(H^a\),
where \(a=v_H(f_0)\).  At each permitted finite tame
specialization, the limiting numerator contains all the zeros
of the limiting denominator with their multiplicities.  This
follows from Lemma~\ref{models:poles}: on a smaller closed disk about a zero,
the initially integral numerator and denominator converge
uniformly, and Weierstrass preparation preserves the number of
zeros with multiplicity.  Thus the specialized division
remainder is zero.  Finite-character detection makes the
remainder zero in \(R\).  There are therefore no remaining
height-one poles.  Since \(R\) is factorial, \(U\in R\), proving
the proposition.
\end{proof}

\subsection{A tame height in analytic rank one}

The preceding argument established integrality without identifying
the central coordinate.  In rank one that identification requires
the following exact rational equality.  Return to a fixed support
\(S\), with no moving primes.  Choose its square order to have
prime support exactly \(S_f\).  Let \(z_b\) be the smoothed
class at the trivial character and put
\[
 d_0=\Delta(0)\prod_{q\in S_f}P_q(1),\qquad
 P_q(1)=L_q(E,1)^{-1}.
\]
If the analytic rank is one, \eqref{eq:source-2} makes the
central dual exponential zero.  At non-\(p\) places rational
elliptic cohomology is zero, by local duality and the Euler
characteristic.  Thus \(z_b\) lies in the rational compact
Selmer line.  Finite \(\Sha\) identifies this line with
\(E(\Q)\otimes\Q_p\), so \(\log_\omega z_b\) means its
local elliptic logarithm.  A nontorsion rational point has
nonzero local logarithm, since the kernel on local points is
torsion.

\begin{proposition}[Exact logarithmic comparison]\label{single:logarithm}
If \(\ord_{s=1}L(E,s)=1\) and \(P\in E(\Q)\) is nontorsion,
then
\begin{equation}
 \frac{\log_\omega z_b}{(\log_\omega P)^2}
       =\pm d_0\frac{L'(E,1)}{\Omega_0H(P)}.
       \tag{5}\label{eq:source-5}
\end{equation}
The right-hand quotient is rational, and is interpreted in
\(\Q_p\).  The height is exactly the height \(H\) fixed in the
introduction.
\end{proposition}

Choose \(K_b=\Q(\sqrt{-D})\), with \(-D<-4\) an odd fundamental
discriminant, so that all primes of \(2Np\) and the other fixed
data split, and
\(L(f\otimes\varepsilon,1)\ne0\), where
\(\varepsilon=\varepsilon_{K_b}\).
Lemma~\ref{classical:auxiliary-fields} applies because \(w(E)=-1\).  Choose
an integral polynomial \(\mathcal A\) in Hecke operators away
from \(NDp\) that kills every other level-\(N\) eigensystem and
the Eisenstein types unramified away from \(N\) with conductor
dividing \(N\), including the weight-two trivial pair, while
\(\mathcal A_f\ne0\).
Also choose a good inert prime \(\ell\nmid6NDp\), away from
the Hecke primes occurring in \(\mathcal A\), with
\(a_f(\ell)\ne0\).  Multiplicity one gives the first choice;
open image and Chebotarev give the second.  The finitely many
Hecke coefficient tests have indices \(n\) prime to \(ND\),
with odd \(\ell\)-valuation.

A horizontal sequence consists of primes \(r_i\), integers
\(m_i\to\infty\), and surjections
\[
 \lambda_i:(\Z/r_i\Z)^\times\longrightarrow\Z/p^{m_i}\Z.
\]
The primes \(r_i\) avoid the level, all fixed Hecke primes, and
all the finitely many tested indices.
We require \(r_i\equiv1\pmod{ND}\), \(r_i\to1\) in \(\Z_p\),
and splitting in the Hilbert class field of \(K_b\).  The
characters kill \(p\), the fixed support, smoothing integers,
\(2,3,D\), all tested indices, and the finite model and lift
exceptions introduced below.  Their limiting Tate Frobenius
has determinant one and trace \(a_*\ne\pm2\), with eigenvalues
not roots of unity.  These characters are automatically even
because \(p\) is odd.  Viewed idelically, \(\lambda_i\) has
uniformizer value \(\lambda_i(q)\) at \(q\ne r_i\), negative
residue character at \(r_i\), and zero uniformizer value there.
Over an extension it is composed with the norm.

\begin{lemma}[A nonzero height direction]\label{single:height-direction}
The horizontal sequence can be chosen to define bilinear tame
heights whose rational limits satisfy
\[
 h_\lambda^{K_b}(P_1,P_2)=2h_\lambda^\Q(P_1,P_2)
 \quad(P_1,P_2\in E(\Q)),\qquad
 h_\lambda^\Q(P,P)\ne0.
\]
These heights are compatible with trace in the second argument.
\end{lemma}

\begin{proof}
Let \(G_y\) be the extension of \(E\) by \(\mathbb G_m\) defined
by the Poincare biextension at \(y=\operatorname{pol}(P)\).
Choose a rational lift \(B_P\) of \(P\).  The fiber is a line
torsor, so this lift exists.  Splitting its underlying Tate
modules gives an upper-right Galois block \(\beta(g)g_T\);
\(\beta\) is the Kummer cocycle of \(y\), up to the fixed
polarization sign.  At finite order this follows by identifying
division points of \(y\) with splittings of the extension and
using the Weil pairing to compare two splittings.  Write the
Kummer cocycle of \(B_P\) as \((k,x)\).  Then
\[
                         dk=-\beta\cup x.
\]
Include all ramification and nonintegrality primes of these
data in the finite killed list.

Adjoin to a fixed normal splitting field all cyclotomic
\(p\)-power roots and all \(p\)-power radicals of the killed
primes.  Its relative Galois group has bounded derived length,
so the Tate image over this field still contains a deep
\(\operatorname{SL}_2\).  On the additional Tate kernel the
values of \(x\) span the plane.  Indeed nonzero Kummer classes
remain nonzero over a fixed finite extension; a central
homothety in the open Tate image detects them on its kernel.
The further radical quotient has scalar cyclotomic conjugation,
so cannot support the irreducible Tate plane as a quotient.
Choose an element \(g\) fixing the splitting and radical data,
with determinant-one Tate action of the required trace, and
\begin{equation}\label{single:height-cocycle}
       k(g)-\beta(g)g_T(g_T-1)^{-1}x(g)\ne0.
\end{equation}
For existence, on the joint kernel the commutator of two
elements has \(x=\beta=0\) and \(k\)-value
\(\pm2e(x(h_1),x(h_2))\), where \(e\) denotes the Weil form.
Full span supplies a nonzero value.
Multiplying the chosen \(g\) by such a commutator changes the
left side of \eqref{single:height-cocycle} without changing its
Tate action.  Chebotarev and compatible power-residue bases now
give the asserted \(r_i,\lambda_i\).

We use the biextension construction of \cite[Section~6,
``Tame heights and trace compatibility'']{B}, verifying its
\(p\)-power version here.  Over an extension
containing a ring class field of positive \(r_i\)-power
conductor, the normed character is unramified also above
\(r_i\).  Indeed the split order class sequence identifies the
two residue-unit factors modulo the diagonal and the global
units \(\{\pm1\}\).  For units at one specified split place,
the other component is one.  The local norm subgroup therefore
has residue in \(\{\pm1\}\), which the even character kills.
Sum the local character valuations of a rational
biextension lift relative to its rigidified integral fibers.
Reciprocity gives independence of the lift, and the two
biextension laws give bilinearity.  On a Jacobian this is the
Deligne-pairing intersection formula for disjoint generic
supports, with the common orientation sign.

To define the same height downstairs, note that
\(\#E(\F_{r_i})\to2-a_*\ne0\) \(p\)-adically.  One fixed
\(p^b\) therefore removes every finite \(p\)-primary reduction
part.  After multiplying the first argument by \(p^b\), divide
it by \(p^{m_i}\) in the uniquely \(p\)-divisible local factor.
For lifts \(\mathcal B,\mathcal B'\) before and after division,
apply the local character to
\[
                    \mathcal B/([p^{m_i}]_1\mathcal B').
\]
Changing \(\mathcal B'\) changes this scalar by a
\(p^{m_i}\)-th power.  The two laws give bilinearity; comparison
with integral lifts gives the unramified construction after
extension.  Norming lifts in the second argument proves trace
compatibility.  Passing to limits and dividing by \(p^b\)
defines the stated rational heights.

For the fixed lift \(B_P\), all places except \(r_i\) contribute
zero.  Put \(w_i=(g_{i,T}-1)^{-1}x(g_i)\), with \(p^bw_i\)
integral after enlarging the fixed \(b\).  Correcting division
points by this coboundary leaves the fiber cocycle
\(p^b(k-\beta\cup w_i)\).  Its Frobenius value is the power
residue character of the scalar ratio above.  Thus the limiting
height is, up to sign, \eqref{single:height-cocycle}, and is
nonzero.  Trace over \(K_b/\Q\) contributes its degree two.
\end{proof}

Use the square classes with the additional prime \(r_i\), and
write \(\chi_b=(1+u)^{\lambda_i}\).  Since \(\lambda_i(p)=0\),
the local coefficients at \(p\) are split.  Let \(Z_i(u)\) be
the group polynomial of dual-exponential coordinates in
\(\omega\).  Local splitting makes every group coefficient
pass through the same fixed map
\(\exp^*:H^1(\Q_p,T)\to\Q_p\omega\).  Its image has a
fixed denominator lattice, so the denominators of \(Z_i\)
are bounded independently of \(i\).  Moreover \(Z_i(0)=0\).
For nontrivial characters,
\[
              Z_i=\pm d_0 L(f,\overline{\chi_b},1)/\Omega_0.
\]
All these characters are primitive at \(r_i\); there is no
omitted \(r_i\)-factor in this formula.  Write \(Z'\) for the
linear coefficient of the bounded-series limit.

\begin{lemma}[Cup comparison]\label{single:cup}
For the sequence of Lemma~\ref{single:height-direction},
\begin{equation}
 Z'\log_\omega P
  =\pm(2-a_*)\frac{\log_\omega z_b}{\log_\omega P}
                        h_\lambda^\Q(P,P).
 \tag{6}\label{eq:source-6}
\end{equation}
\end{lemma}

\begin{proof}
At stage \(i\), write a cocycle modulo \((p^{m_i},u^2)\) as
\(z_0+uz_1\).  Then \(dz_1=-\lambda_i\cup z_0\).
The good norm relation and one fixed clearing factor permit
the adjustment
\[
 z_0=\alpha_i\beta,\qquad
 \alpha_i=\pm\left(1-\frac{a_f(r_i)}{r_i}+\frac1{r_i}\right)
                              \frac{\log_\omega z_b}{\log_\omega P}.
\]
The clearing is uniform because the degree-one torsion has
bounded exponent.  The scalar two-cocycle
\(z_1\cup x+\alpha_i\lambda_i\cup k\) is closed.  Its local
invariant at \(p\) tends to \(Z'\log_\omega P\), by the
logarithm--dual-exponential pairing.  At the other fixed places
the character is killed and the rational local class is a
coboundary.  At \(r_i\), use \(dw_i=x\) to obtain
\(\alpha_i\lambda_i\cup(k-\beta\cup w_i)\).  Its invariant is
\(\pm\alpha_i h_i^\Q(P,P)\), by the scalar-ratio computation
in Lemma~\ref{single:height-direction}.  The sum of local
invariants is zero.  Passing to the limit, the Euler factor
tends to \(2-a_*\), giving \eqref{eq:source-6}.
\end{proof}

\subsection{The exact spectral comparison}

Fix a rational modular parametrization \(\phi:X_0(N)\to E\),
of degree \(\delta_\phi\).  Let \(P_X\in J_0(N)(K_b)\) be the
single oriented Hilbert-class trace based at \(\infty\), and
put \(Q_X=\phi_*P_X\).  With unnormalized Petersson integral set
\[
                        C_0=\frac{\sqrt D}{8\pi^2(f,f)_N}.
\]
Choose a real primitive period \(\Omega_{\rm tw}\) for the
negative twist.  The period relations are
\[
 C_0\Omega_0\Omega_{\rm tw}\in\Q^\times,
 \qquad L(f\otimes\varepsilon,1)/\Omega_{\rm tw}\in\Q^\times.
\]
The second is modular-symbol rationality.  For the first, if
\(\phi^*\omega=c_\phi f\,dq/q\), the complex area of \(E\)
is \(8\pi^2c_\phi^2(f,f)_N/\delta_\phi\); the twist period is a
rational multiple of the anti-invariant period divided by
\(\sqrt D\).  This also specifies the rational normalization
in which the following identities are embedded in \(\Q_p\).

\begin{lemma}[Spectral-height identity]\label{single:spectral}
For the horizontal sequence just constructed,
\begin{equation}
 C_0\Omega_0L(f\otimes\varepsilon,1)\frac{Z'}{d_0}
       =\pm(2-a_*)\frac{h_\lambda^{K_b}(Q_X,Q_X)}{\delta_\phi}.
       \tag{7}\label{eq:source-7}
\end{equation}
\end{lemma}

\begin{proof}
The proof evaluates the same finite differences of kernel
coefficients in two ways.  Projection onto the \(f\)-block
expresses them through \(Z'\); their intersection interpretation
expresses them through \(h_\lambda^{K_b}(Q_X,Q_X)\).  Matching
the two evaluations, with their constants, gives the identity.
We first specify the finite-character formulas and then prove
their interpolation in the present \(p\)-power direction.
They are the complex coefficient and geometric intersection
identities in \cite[Section 6, ``The traced coefficient,''
``Fricke transformation and unfolding,'' and
``Intersection realization'']{B}.
Their hypotheses are an odd fundamental discriminant \(-D<-4\),
split Heegner level \(N\), a prime \(r\nmid ND\) split in the
Hilbert class field with \(r\equiv1\pmod{ND}\), an even
prime-conductor character \(\chi\) with \(\chi^2\ne1\), and
indices \(nr^h\) with \((n,ND)=1\) and odd valuation at the
chosen inert prime \(\ell\).  None of these identities has a
coefficient-prime hypothesis.

Here are their precise normalizations.  Choose integral ideal
representatives \(C\) of \(\Pic(\mathcal O_{K_b})\), prime to
\(D\), and put \(a_C=\mathrm NC\).
Let \(e(x)=\exp(2\pi i x)\) and
\[
 \theta(z)=\frac12\sum_C\sum_{x\in C}e(\mathrm Nx\,z/a_C).
\]
Its positive coefficient is
\(\rho(n)=\sum_{d\mid n}\varepsilon(d)\).  If \(\theta_\chi\)
is its Fourier twist, define
\[
 E_{\chi,s}(z)=\frac12
  \sum_{\substack{(c_1,d_1)\ne(0,0)\\ND\mid c_1}}
    \varepsilon(d_1)\bar\chi^2(c_1)
       \frac{\operatorname{Im}(z)^s}
            {(c_1z+d_1)|c_1z+d_1|^{2s}},
 \qquad
 I_\chi=\left.\frac{\sqrt D}{2\pi}
    \Tr_{NDr^2/Nr^2}(\theta_\chi E_{\chi,s})\right|_{s=0}.
\]
Characters are extended by zero at nonunits; the trace is
unnormalized.  The continuation is holomorphic.  Poisson
summation pairs the odd row weights, eliminates the
nonholomorphic modes, and permits pairing with cusp tests.

For every tested index \(n\) and \(h\ge1\), regard the
coefficient at \(nr_i^h\) as a polynomial
\(D_{i,h}(n;u)\) in
\(\Z_p[u]/((1+u)^{p^{m_i}}-1)\): in the following sum
\(r=r_i\) and \(\chi(x)=(1+u)^{\lambda_i(x)}\).
\begin{equation}\label{single:kernel-coefficient}
 D_{i,h}(n;u)=\frac12
 \sum_{\substack{C,\ y\in C,\ b>0\\
       a=\mathrm Ny/a_C>0\\ a+Nb=nr^hD}}
 \sum_{\substack{kl=b\\\gcd(k,l,D)=1}}
 W(k)\chi(a/D)\bar\chi^2(Nk),
\end{equation}
where
\[
 W(k)=\prod_{\substack{q\mid D\\q\mid k}}(l/q)
            \prod_{\substack{q\mid D\\q\nmid k}}(-a_CNk/q).
\]
Write \(\epsilon_q(x)=(x/q)\).  The local partial Fourier
transform at \(q\mid D\) has scalar
\(\kappa_q=q^{-1}\tau_q\epsilon_q(-a_CD/q)\), with
\(\tau_q=\sum_{x\bmod q}\epsilon_q(x)e(x/q)\).
Multiplying the partial traces
and the Poisson factor \(-2\pi i/D\) leaves common scalar one,
because the primitive Gauss sum of \(\varepsilon\) is
\(i\sqrt D\).  This explains in particular the factor \(1/2\)
in \eqref{single:kernel-coefficient}.  For a trace matrix
\(\gamma=\left(\begin{smallmatrix}a_\gamma&b_\gamma\\
c_\gamma&d_\gamma\end{smallmatrix}\right)\in\Gamma_0(Nr^2)\),
translation of the theta twist multiplies its weight by
\(\chi(a_\gamma^{-2})\), canceled by the Eisenstein row
weight \(\chi^2(a_\gamma)\).  Thus these
finite-character computations apply to our characters as soon
as \(\chi^2\ne1\), which holds for every nontrivial
\(p\)-power character.

Use the same marked Artin limits as before to define
\(D_h(n;u)=\lim_i D_{i,h}(n;u)\) and
\(D'_h(n)=[u]D_h(n;u)\).  Only terms with \(a\) a unit at
\(r\) occur in \eqref{single:kernel-coefficient}.  Since
\(r\mid nr^h\), this is equivalent to \(k,l\) being units
there.  The identity \(W(l)=-W(k)\) cancels the
divisor-independent terms in the derivative of the character
weight.  Hence \(D'_h(n)\) is the limit of the masked sum
with weight \(-2\lambda_i(k)\).

For \(h\ge0\), define \(B_{i,h}(n)\in\Z/p^{m_i}\Z\) by
the same sum with the mask removed and character weight
replaced by \(-2\lambda_i(k/r_i^{v_{r_i}(k)})\).  Write
\(B_h(n)\) for its Artin limit, put \(B_h(n)=0\) for
\(h<0\), and set \(C_h(n)=B_h(n)-B_{h-2}(n)\).
Stripping a common power \(r^j\) from \(a,b\) shows that,
for \(h\ge2\),
\begin{equation}\label{single:mask}
 C_{h+1}-2C_h+C_{h-1}=D'_{h+1}+2D'_h+D'_{h-1}.
\end{equation}
For completeness, when the primitive index is positive the
\(C_h\)-multiplicities are \(1,4,8,12,\ldots\); their second
differences are \(1,2,1\).  At primitive index zero they are
affine and the second difference is zero.  The off-\(r\)
weights do not change, since the primes over \(r\) are principal
and \(r\equiv1\pmod D\).  This proves the identity for
\(\Z/p^{m_i}\)-valued weights and then on their limits, with
no alteration of constants.

The symbols \(D'_h,B_h,C_h\) will also denote these coefficient
functions of \(n\).  The Hecke polynomial \(\mathcal A\)
acts on them by the weight-two coefficient rule
\((T_s a)(n)=a(sn)+s a(n/s)\), with the second term zero
unless \(s\mid n\).  All its primes are away from \(NDr_i\).
The finite set of tested indices was chosen to contain the
indices needed to evaluate this polynomial at \(n=\ell\).

We next project its holomorphic side.  At level \(Nr^2\) use
\(f_j=r^jf(r^jz)\), \(0\le j\le2\), and write
\(\tau(\chi)=\sum_{x\bmod r}\chi(x)e(x/r)\).
The finite-character Fricke and unfolding identities are
\begin{equation}\label{single:old-pairings}
 \frac{(I_\chi,f_2)_{Nr^2}}{(f,f)_{Nr^2}}
     =\frac{C_0}{r+1}\frac{\tau(\chi)}{\tau(\bar\chi)}
                L(f,\bar\chi,1)L(f\otimes\varepsilon,\bar\chi,1),
 \qquad (I_\chi,f_1)_{Nr^2}=0.
\end{equation}
The exact constants arise from the kernel factor
\(\sqrt D/(2\pi)\), the Fricke factor \(r\), the Mellin integral
\(1/(4\pi)\), and the level index \(r(r+1)\).
At primes of \(N\), use the primitive degree-at-most-one
polynomial and \(\varepsilon(q)=1\); at primes of \(D\) the
twisted local factor is one; at \(r\) both twisted factors
are one.  The total nebentypus is trivial.  The conductor-\(r\)
new cuspidal block is therefore Steinberg, with \(U_r=\pm1\),
and its complementary raises have eigenvalue zero.  The
conductor-\(r^2\) cuspidal block has \(U_r=0\).  In a
ramified Eisenstein pair both characters are ramified at \(r\),
so its \(U_r\)-eigenvalue is also zero.
Thus \(U_r^h(U_r^2-1)\) kills these blocks for sufficiently
large \(h\).  The remaining old blocks come from level \(N\);
the fixed operator \(\mathcal A\) kills all of them except
the \(f\)-block.

Write \(a_h=a_f(r^h)\) and \(a=a_f(r)\).  The normalized
Gram matrix of \(f_0,f_1,f_2\) has entries \(g_{|i-j|}\), where
\[
 g_0=1,\qquad g_1=\frac a{r+1},\qquad rg_2=ag_1-1.
\]
The coefficient vector at index \(nr^h\), after removing its
\(n\)-factor, is \((a_h,ra_{h-1},r^2a_{h-2})\).  The
coefficient of the Gram row for \(f_2\) is
\[
 s_h=\frac{r^2a_{h-2}-g_1ra_{h-1}}{1-g_1^2}.
\]
Its denominators have bounded \(p\)-valuation because
\(a_*\ne\pm2\).  Set
\(\gamma_h=\lim_i(a_f(r_i^h)-a_f(r_i^{h-2}))\).  The
Hecke recurrence gives
\[
 \lim_i\frac{s_h}{r_i+1}
       =\frac{\gamma_h}{(a_*-2)(a_*+2)}.
\]

All factors in \eqref{single:old-pairings} have bounded group
polynomial realizations.  The Gauss ratio is
\(\tau(\chi)^2/r\); its coefficients are integral and fixed
by \(p\)-Frobenius because \(\lambda_i(p)=0\), and its
limiting augmentation is one.  The first \(L\)-value factor
is \(Z_i/d_0\).  For the companion, use the even additive
modular-symbol sum multiplied by \(\tau(\bar\chi)/r\).
Its denominators are bounded by the fixed relative period
lattices.  Before Gauss multiplication its augmentation is
\[
 r\sum_{r\mid j}\frac{a_{f\otimes\varepsilon}(j)}j
    -L(f\otimes\varepsilon,1)
      =(a_f(r)-2)L(f\otimes\varepsilon,1),
\]
by the Hecke recurrence and Mellin continuation.  The trivial
Gauss sum is \(-1\), so its limiting normalized augmentation
is \((2-a_*)L(f\otimes\varepsilon,1)/\Omega_{\rm tw}\).
Multiply the finite-character identity by \((1+u)^2-1\) to
remove the excluded character values, pass to bounded-series
limits, and cancel the nonzero multiplier.  Since \(Z_i(0)=0\),
only that companion augmentation contributes to the derivative.
We obtain
\begin{equation}\label{single:projected-derivative}
 \begin{split}
 [\mathcal A(D'_{h+2}-D'_h)](\ell)
   ={}&\pm\mathcal A_fa_f(\ell)
       \frac{\gamma_{h+2}-\gamma_h}{(a_*-2)(a_*+2)}\\
      &\hspace{4mm}\cdot C_0\Omega_0\frac{Z'}{d_0}
                     (2-a_*)L(f\otimes\varepsilon,1).
 \end{split}
\end{equation}
This interpolation uses bounded denominators, not a
coefficient-prime restriction; in particular it applies when
\(p=3\).

For the geometric side, the coefficient-prime-independent
intersection identity cited above applies at
\(q\nmid6NDnr^h\).  Its local divisor sum, with weight
\(-2v_q(k)\) and the original factor \(1/2\), is zero unless
\(q\) is inert, \(d_q=v_q(b)\) is odd, and
\(\varepsilon_{K_b,s}(-bNa_C)=1\) for all \(s\mid D\), where
\(\varepsilon_{K_b,s}\) is the local quadratic norm character
of \(K_b/\Q\) at \(s\), distinguished from the residue symbols
\(\epsilon_q\) above.
When these conditions hold it equals
\[
 -\frac{1+d_q}{2}\rho(b/q)
                   2^{\#\{s\mid D:s\mid b\}}.
\]
Its negative is the intersection of the oriented ordered
Hilbert-class sums, based respectively at \(\infty\) and
\(0\), with the second acted on by \(T_{nr^h}\), weighted by
residue degree over \(q\).  The generic supports are disjoint
by the odd inert \(\ell\)-valuation.  The remaining fixed
primes have character zero.  This yields all the required
height contributions except those at \(r_i\).

The difference \(C_h=B_h-B_{h-2}\) keeps cyclic paths in the
second Hecke divisor; its second difference eliminates the two
continuing maximal-order paths.  Every other orbit is a trace
from an extension containing a positive-\(r\)-conductor ring
class field.  After pulling up the first divisor, intersections
at \(r\) are zero.  Indeed ordinary CM reduction gives linear
or conjugate-semilinear maps over the unramified quadratic
\(\ell\)-adic algebra, whose norm degrees have even
\(\ell\)-valuation; they cannot produce the prescribed odd
index.  Cusps remain disjoint.  The missing contributions may
therefore be added as zeros.  Lemma~\ref{single:height-direction}
now identifies the resulting intersections with full trace
heights: the normed characters are unramified over exactly the
extensions just obtained.

Applying \(\mathcal A\) at \(\ell\) inserts the projector
\(\phi^*\phi_*/\delta_\phi\), with scalar
\(\mathcal A_fa_f(\ell)\).  The Hecke second difference acts
on the second point with eigenvalue \(\gamma_h\).  Thus the
geometric side of \eqref{single:mask}, after the additional
\(h\mapsto h+2\) difference, is
\[
 \pm\mathcal A_fa_f(\ell)(a_*-2)(\gamma_{h+2}-\gamma_h)
                    \frac{h_\lambda^{K_b}(Q_X,Q_X)}{\delta_\phi}.
\]
All projector denominators and the fixed first-argument
multiplier \(p^b\) cancel in these rational identities.
Cusp torsion has zero rational height.  On the holomorphic
side, the coefficients \(1,2,1\) in \eqref{single:mask}
insert \(a_*+2\) into \eqref{single:projected-derivative}.
Cancel \(\mathcal A_fa_f(\ell)\) and
\(\gamma_{h+2}-\gamma_h\).  The latter is nonzero for some
sufficiently large \(h\), since the limiting eigenvalues are
distinct and not roots of unity.  The remaining scalar
\((a_*+2)(2-a_*)/((a_*-2)(a_*+2))=-1\) gives
\eqref{eq:source-7}.
\end{proof}

\begin{proof}[Proof of Proposition~\ref{single:logarithm}]
The companion has analytic rank zero.  Hence
\(E(K_b)\otimes\Q=E(\Q)\otimes\Q\), and we may write
\(Q_X=t_PP\) modulo torsion, with \(t_P\in\Q\).  The exact
split Gross--Zagier formula of Theorem~\ref{classical:gross-zagier}, for
this single Hilbert-class trace, is
\begin{equation}\label{single:gross-zagier}
 C_0L'(f,1)L(f\otimes\varepsilon,1)
                     =\frac{2t_P^2H(P)}{\delta_\phi}.
\end{equation}
In particular \(t_P\ne0\).  Trace compatibility gives
\(h_\lambda^{K_b}(Q_X,Q_X)=2t_P^2h_\lambda^\Q(P,P)\).
The same factor two occurs in \eqref{single:gross-zagier}.
Substituting both identities into \eqref{eq:source-7} gives
\[
 Z'=\pm d_0(2-a_*)\frac{L'(f,1)}{\Omega_0H(P)}
                                      h_\lambda^\Q(P,P).
\]
Compare with \eqref{eq:source-6}.  The factors
\(2-a_*\), \(h_\lambda^\Q(P,P)\), and \(\log_\omega P\)
are nonzero.  Their cancellation proves
\eqref{eq:source-5}, with no unit ambiguity other than the
common orientation sign.
\end{proof}

\subsection{The central integral calculation}

\begin{proposition}[Central coordinate]\label{single:center}
The element \(U\) of Proposition~\ref{single:integral} satisfies
\begin{equation}
                         X(E)=v_p(U(0))\ge0.
                         \tag{8}\label{eq:source-8}
\end{equation}
This holds at every odd coefficient prime, for all reduction
types and all residual representations of the non-CM curve.
\end{proposition}

\begin{proof}
Compute first at an actual stage with trivial characters and
support \(S'=S\cup\{r_1,\ldots,r_k\}\).  Put
\[
 O^j=H^j(G_{\Q,S'},T),\quad
 s_0=v_p(\#\Sha(E/\Q)),\quad
 \tau_g=v_p(\#E(\Q)_{\rm tors}),\quad
 \tau_v=v_p(\#E(\Q_v)[p^\infty]).
\]
Let \(p^l\Z_p=\log_\omega E(\Q_p)\).  The singular quotient
\(J_p\) of local degree-one cohomology by compact Kummer is a
free line, perfectly dual to compact Kummer modulo torsion.
At non-\(p\) places Kummer fills compact degree one.  Integral
Poitou--Tate therefore gives
\begin{equation}\label{single:pt}
 \begin{split}
 O^1\longrightarrow J_p\longrightarrow\Sel_{p^\infty}(E/\Q)^\vee
 \longrightarrow O^2\longrightarrow
 \bigoplus_{v\in S'_f}H^2(\Q_v,T)
 \longrightarrow E(\Q)[p^\infty]^\vee\longrightarrow0.
 \end{split}
\end{equation}
Real Tate terms vanish because \(p\) is odd.  The local
degree-two groups have lengths \(\tau_v\).  Also \(O^0=0\),
\(\rank O^1=1\), and \(\length O^1_{\rm tors}=\tau_g\).

In rank zero, put
\(j=\length(J_p/\im O^1)\).  Equation~\eqref{single:pt} gives
\[
              \length O^2=s_0-j+\sum_{v\in S'_f}\tau_v-\tau_g.
\]
If \(z_i^\circ\) is the central stage class and
\(\exp^*z_i^\circ=\alpha_i\omega\), logarithm--exponential
duality gives its valuation in \(J_p\) as \(v_p(\alpha_i)+l\).
Its valuation in \(O^1/O^1_{\rm tors}\) is therefore
\(v_p(\alpha_i)+l-j\).  Adding degree-one torsion and
subtracting degree-two torsion gives raw determinant order
\begin{equation}\label{single:raw-zero}
          v_p(\alpha_i)+l+2\tau_g-s_0-\sum_{v\in S'_f}\tau_v.
\end{equation}

In rank one, let \(P\) generate the full free Mordell--Weil
group modulo torsion.  The map from \(J_p\) to the dual of
the divisible global point subgroup is injective, with
cokernel of length \(v_p(\log_\omega P)-l\).  Consequently
\[
 \length O^2=s_0+v_p(\log_\omega P)-l
                                  +\sum_{v\in S'_f}\tau_v-\tau_g.
\]
Finite \(\Sha\) identifies \(O^1\) with compact global points.
The free coordinate of the class has valuation
\(v_p(\log_\omega z_i^\circ)-v_p(\log_\omega P)\).
The same signed torsion sum therefore gives
\begin{equation}\label{single:raw-one}
 v_p\left(\frac{\log_\omega z_i^\circ}
                        {(\log_\omega P)^2}\right)
                   +l+2\tau_g-s_0-\sum_{v\in S'_f}\tau_v.
\end{equation}
The logarithm is nonzero by \eqref{eq:source-5}.  When restoring
the moving omissions, the good norms multiply the fixed class
by \(\prod_r(1-a_f(r)/r+1/r)\), with exactly these factors.

Lemma~\ref{models:haar} gives the local Haar identities for
the minimal differential:
\[
 \tau_q=v_p(c_qP_q(1))\quad(q\ne p),\qquad
 \tau_p-l=v_p(c_pP_p(1)).
\]
They include all reduction types: the nonsingular reduction
measure is its point count divided by the residue cardinality,
and the formal logarithm accounts for the finite torsion kernel
at \(p\).  Substitute \eqref{eq:source-2} into
\eqref{single:raw-zero}, or \eqref{eq:source-5} into
\eqref{single:raw-one}.  Every omitted Euler factor cancels
against its Haar factor.  The result in either case is
\[
                         X(E)+v_p(\Delta(0)).
\]
In rank one the use of a generator of the full free group makes
\(H(P)=\Reg_E\).  The ratio \(\Omega_E/\Omega_0\) is a
power of two and has zero \(p\)-valuation.  Thus neither a
Mordell--Weil index nor a real-component factor remains.

Finally pass from the stage calculation to the central value
of the limit.  The moving local torsion lengths are bounded,
because their Frobenius traces tend to values different from
two.  The displayed cohomology formulas bound the remaining
torsion lengths.  Hence determinant pivots of the same ranks
survive the limit by Remark~\ref{models:central-pivots}.
The raw valuation is unchanged.
Each \(D_q/P_q\) has central value one; division by
\(\Delta(0)\) leaves \(v_p(U(0))=X(E)\).
Proposition~\ref{single:integral} now gives the inequality.
\end{proof}
